%!TEX TS-program = xelatex
%!TEX encoding = UTF-8 Unicode

% Three-page CV summary tailored to the UNESCO IGCP "Call for Expressions of
% Interest: Geoscience Funding for Latin America and the Caribbean" (deadline
% 13 Jul 2026). Uses the same Awesome-CV design as the full Spanish/English CVs.
\PassOptionsToPackage{left=5cm,top=2cm,right=1.5cm,bottom=2.5cm,nohead,nofoot}{geometry}
\documentclass[10pt,a4paper]{../awesome-cv}

\fontdir[fonts/]

\geometry{left=1.4cm,top=1.3cm,right=1.4cm,bottom=1.3cm,nohead,nofoot}
\colorlet{awesome}{awesome-red}
\setbool{acvSectionColorHighlight}{true}
\renewcommand{\acvHeaderSocialSep}{\quad\textbar\quad}

\usepackage[english]{babel}
\usepackage{enumitem}
\usepackage{multicol}
\usepackage{needspace}

\photo{../profile.png}
\name{Jeremías}{Likerman}
\position{Adjunct Researcher, CONICET{\enskip\cdotp\enskip}Associate Professor, University of Buenos Aires}
\address{Institute of Andean Studies ``Don Pablo Groeber'' (IDEAN), Argentina}
\email{jlikerman@uba.ar}
\googlescholar{HGA3sQQAAAAJ}{Google Scholar}
\extrainfo{\href{https://orcid.org/0000-0003-3844-2085}{ORCID 0000-0003-3844-2085}}

\makeatletter
\renewcommand*{\sectionstyle}[1]{{\fontsize{12.5pt}{15pt}\bodyfont\bfseries\color{text}\@sectioncolor #1}}
\renewcommand*{\subsectionstyle}[1]{{\fontsize{9.6pt}{11.5pt}\bodyfont\bfseries\scshape\textcolor{graytext}{#1}}}
\renewcommand*{\paragraphstyle}{\fontsize{9pt}{11.8pt}\bodyfontlight\upshape\color{text}}
\renewcommand*{\entrytitlestyle}[1]{{\fontsize{9.6pt}{11.5pt}\bodyfont\bfseries\color{darktext} #1}}
\renewcommand*{\entrypositionstyle}[1]{{\fontsize{8.4pt}{10pt}\bodyfont\scshape\color{graytext} #1}}
\renewcommand*{\entrydatestyle}[1]{{\fontsize{8.4pt}{10pt}\bodyfontlight\slshape\color{graytext} #1}}
\renewcommand*{\entrylocationstyle}[1]{{\fontsize{8.6pt}{10.3pt}\bodyfontlight\slshape\color{awesome} #1}}
\renewcommand*{\descriptionstyle}[1]{{\fontsize{8.8pt}{11.2pt}\bodyfontlight\upshape\color{text} #1}}
\renewcommand*{\honorpositionstyle}[1]{{\fontsize{8.8pt}{11pt}\bodyfont\mdseries\color{darktext} #1}}
\renewcommand{\cvsection}[1]{%
  \par\vspace{3mm}%
  \sectionstyle{#1}%
  \phantomsection%
  \color{gray}\vhrulefill{0.9pt}%
  \par\nobreak\vspace{1.8mm}%
}
% Local fix: the class wraps consecutive \cventry tables in a `center`
% environment with no \par between them, which both misaligns entries
% left/right and collapses the gap before the next entry's title. Redefine
% \cventries/\cventry to break cleanly and add explicit spacing instead.
\renewenvironment{cventries}{%
  \vspace{\acvSectionContentTopSkip}
}{%
}
\renewcommand*{\cventry}[5]{%
  \par\noindent
  \setlength\tabcolsep{0pt}
  \setlength{\extrarowheight}{0pt}
  \begin{tabular*}{\textwidth}{@{\extracolsep{\fill}} L{\textwidth - 4.5cm} R{4.5cm}}
    \ifempty{#2#3}
      {\entrypositionstyle{#1} & \entrydatestyle{#4} \\}
      {\entrytitlestyle{#2} & \entrylocationstyle{#3} \\
      \entrypositionstyle{#1} & \entrydatestyle{#4} \\}
    \multicolumn{2}{L{\textwidth}}{\descriptionstyle{#5}}
  \end{tabular*}\par
  \vspace{2.0mm}%
}
\makeatother

\setlist[itemize]{leftmargin=4mm,topsep=0.4mm,itemsep=0.4mm,parsep=0mm}
\setlength{\emergencystretch}{1.5em}
\setlength{\footskip}{17pt}

\begin{document}

\makecvheader

\makecvfooter
  {\today}
  {Jeremías Likerman ~~\textperiodcentered~~ CV summary --- UNESCO IGCP LAC call}
  {\thepage}

%-------------------------------------------------------------------------------
\cvsection{Profile}
\begin{cvparagraph}
Geoscientist (PhD, 2015) specializing in structural geology, geomechanics, and multiscale numerical modelling of the Andean orogen, based at CONICET--IDEAN (Argentina), an institution eligible under this call. Fifteen years of Andean field research address two of the six IGCP priority themes: \textbf{Geodynamics} (subduction dynamics, lithospheric deformation, and mantle-plume interactions along the Andean margin) and \textbf{Earth Resources} (natural-fracture prediction and connectivity in the Vaca Muerta unconventional play, and geothermal-potential assessment of volcanic systems). A sustained record of teaching, professional training, thesis supervision, and prior direct collaboration with researchers in Chile supports the regional-cooperation and capacity-building priorities of this call.
\end{cvparagraph}

%-------------------------------------------------------------------------------
\cvsection{Education}
\begin{cventries}
\cventry{University of Buenos Aires}{PhD in Geological Sciences}{Argentina}{2010--2015}{\begin{cvitems}\item{Thesis: \emph{Prediction of fracturing on irregular geological surfaces using static methods}. Supervisor: Dr. Ernesto Cristallini.}\end{cvitems}}
\cventry{CONICET and Universitat Politècnica de Catalunya}{Postdoctoral Fellowship}{Argentina/Spain}{2015--2017}{\begin{cvitems}\item{Thermal evolution of oceanic lithosphere. Supervisors: Dr. Ernesto Cristallini and Dr. Sergio Zlotnik.}\end{cvitems}}
\cventry{University of Buenos Aires}{BSc in Geological Sciences}{Argentina}{2005--2010}{\begin{cvitems}\item{Thesis: \emph{Geology and structure of the Tres Cruces compound anticline, Jujuy, Argentina}.}\end{cvitems}}
\end{cventries}

%-------------------------------------------------------------------------------
\cvsection{Current Appointments}
\begin{cventries}
\cventry{CONICET--IDEAN}{Adjunct Researcher (Investigador Adjunto)}{Argentina}{2018--present}{\begin{cvitems}\item{Fracture prediction, geodynamic modelling, and geothermal-potential modelling. Promoted to Investigador Adjunto, Nov. 2023.}\end{cvitems}}
\cventry{Department of Geological Sciences, University of Buenos Aires}{Associate Professor}{Argentina}{May 2026--present}{\begin{cvitems}\item{Courses: Geostatistics and Geoinformatics.}\end{cvitems}}
\cventry{IDEAN, CONICET--UBA}{Institutional Coordinator}{Argentina}{Sep 2025--present}{\begin{cvitems}\item{Institute council member since April 2024 --- direct experience in research-institute governance and administration.}\end{cvitems}}
\end{cventries}

%-------------------------------------------------------------------------------
\Needspace{15\baselineskip}
\cvsection{Research Projects (Selected) --- IGCP-Aligned}
\begingroup
\renewcommand{\arraystretch}{1.05}
\begin{cvhonors}
\cvhonor{Tectonic-structural analysis through analogue and numerical modelling: contributions to energy resources.}{PIP 11220250100817CO, CONICET}{Argentina}{2026--28}
\cvhonor{Geothermal potential of volcanic zones \textbf{(PI)} --- \emph{Earth Resources}.}{AECT-2025-1-0009}{Spain}{2025}
\cvhonor{Structural controls in volcanic systems: intrusions and damage zones linked to hydrocarbon generation \textbf{(Co-PI)} --- \emph{Earth Resources}.}{Williams Foundation}{Argentina}{2025--26}
\cvhonor{Thermomechanical assessment of geothermal potential in volcanic zones \textbf{(PI)} --- \emph{Earth Resources}.}{LINCGL2024}{Spain}{2024--26}
\cvhonor{Geodynamic modelling of subduction zones and mantle plumes \textbf{(PI)} --- \emph{Geodynamics}.}{AECT-2024-2-0027}{Spain}{2024}
\cvhonor{Geodynamics of the Adria microplate, from mantle to surface \textbf{(Co-PI)} --- \emph{Geodynamics}.}{GEO3BCN}{Spain}{2023}
\cvhonor{Geodynamic modelling of subduction zones \textbf{(PI)} --- \emph{Geodynamics}.}{AECT-2022-3-0023}{Spain}{2022}
\cvhonor{Natural fracturing and first-order structures, Cerro Domuyo \textbf{(Co-PI)} --- \emph{Earth Resources}.}{UBACyT}{Argentina}{2018}
\cvhonor{Numerical modelling of the thermal evolution of oceanic lithosphere \textbf{(PI)} --- \emph{Geodynamics}.}{PICT 2015-1226}{Argentina}{2016}
\end{cvhonors}
\endgroup
\descriptionstyle{Five of the nine projects above were led as Principal Investigator, demonstrating a consistent track record of securing and managing competitive, internationally co-funded research grants.}

%-------------------------------------------------------------------------------
\cvsection{Capacity-Building, Teaching \& Regional Engagement}

\cvsubsection{Teaching and professional training}
\begin{cventries}
\cventry{Instructor in charge}{Geology of Unconventional Reservoirs}{Argentina}{2025}{\begin{cvitems}\item{Professional training course on the geological characterization of shale reservoirs.}\end{cvitems}}
\cventry{Field instructor}{Advanced Research Field Trip}{Argentina}{2025}{\begin{cvitems}\item{Organizer and instructor, \emph{Perspectives on Vaca Muerta: subsurface attributes analysed from surface exposures}.}\end{cvitems}}
\cventry{Department of Geological Sciences, UBA}{Geostatistics, Geoinformatics, Geological Surveying, Introduction to Geology}{Argentina}{2013--present}{\begin{cvitems}\item{Continuous university teaching, from Teaching Assistant to Associate Professor.}\end{cvitems}}
\end{cventries}

\cvsubsection{Supervision of early-career researchers}
\descriptionstyle{Supervised 1 postdoctoral researcher (Plotek, B., 2024), 1 doctoral thesis (Correa Luna, C., 2023), and 11 undergraduate theses (2015--present, several ongoing). Currently mentoring 1 CONICET Assistant Researcher (Plotek, B., since Aug. 2025).}

\cvsubsection{Regional and international cooperation}
\begingroup
\fontsize{8.8pt}{11.2pt}\selectfont\bodyfontlight\color{text}
\begin{cvitems}
\item{\textbf{Chile:} multi-year collaboration (2012--2017) with researchers at Universidad de Chile (L. Pinto, R. Charrier) on Andean tectonic inversion, resulting in a co-authored article in the \emph{Journal of South American Earth Sciences} (2017) and a chapter in \emph{Geological Society, London, Special Publications} v.\,399 (2014).}
\item{Participant, SZNet 2025 Chile Field Trip, NSF SZ4D initiative --- ongoing engagement with the regional subduction-zone research network.}
\item{\textbf{Spain (UPC):} postdoctoral research and recurring visiting stays (2015, 2017, 2019, 2022) at the Numerical Computation Laboratory, Universitat Politècnica de Catalunya (BarcelonaTech), under Dr. Sergio Zlotnik --- a decade-long bilateral cooperation.}
\item{Coordinator, CONICET--CSIC (Spain) bi-national cooperation on geothermal modelling (2024--2026) --- direct experience running an internationally co-funded, multi-year cooperation project of comparable scope to this call.}
\item{Incoming Research Fellowship, Kobe University, Japan (JSPS, 2027) --- a further, independently competitive international fellowship, reflecting a sustained track record of securing international research placements.}
\end{cvitems}
\endgroup

%-------------------------------------------------------------------------------
\cvsection{Publications}
\descriptionstyle{Complete list of 19 peer-reviewed journal articles (2013--2026), reverse chronological order.}\par
\vspace{0.8mm}
{\fontsize{7.6pt}{9.6pt}\bodyfontlight\color{text}
\begin{multicols}{2}
\begin{enumerate}
  \setlength{\itemsep}{0.5mm}
  \item Villarroel, M.A., Browning, J., Jara, P.P., Harnett, C.E., Holohan, E.P., Marquardt, C.J., Guzman, M., \& \textbf{Likerman, J.} (2026). The role of compressional structures in the development of collapse calderas: Insights from analogue models. \emph{Journal of Volcanology and Geothermal Research}, 473, 108582.
  \item Gianni, G.M., Gallo, L.C., \textbf{Likerman, J.}, Echaurren, A., Gianni, C.R., \& Faccena, C. (2025). Slab underthrusting is the primary control on flat-slab size. \emph{Science Advances}, 11(28).
  \item Correa-Luna, C., Yagupsky, D.L., \textbf{Likerman, J.}, \& Barcelona, H. (2025). Impact of large-scale structures on fracture network connectivity: Insights into the Vaca Muerta unconventional play, Neuquén Basin, Argentina. \emph{Tectonophysics}, 230640.
  \item Plotek, B., \textbf{Likerman, J.}, \& Cristallini, E. (2024). Geomechanical modeling of fault-propagation folds: A comparative analysis of finite-element and the trishear kinematic model. \emph{Journal of Structural Geology}, 105064.
  \item Navarrete, C., Gianni, G., Tassara, S., Zaffarana, C., \textbf{Likerman, J.}, Márquez, M., Wostbrock, J., Planavsky, N., Tardani, D., \& Perez Frasette, M. (2024). Massive Jurassic slab break-off revealed by a multidisciplinary reappraisal of the Chon Aike silicic large igneous province. \emph{Earth-Science Reviews}, 249, 104651.
  \item Gianni, G.M., \textbf{Likerman, J.}, Navarrete, C.R., Gianni, C.R., \& Zlotnik, S. (2023). Ghost-arc geochemical anomaly at a spreading ridge caused by supersized flat subduction. \emph{Nature Communications}, 14, 2083.
  \item Plotek, B., Heckenbach, E., Brune, S., Cristallini, E., \& \textbf{Likerman, J.} (2022). Kinematics of fault-propagation folding: Analysis of velocity fields in numerical modeling simulations. \emph{Journal of Structural Geology}, 162, 104703.
  \item Correa-Luna, C., Yagupsky, D.L., \textbf{Likerman, J.}, \& Barcelona, H. (2022). Natural fracture characterization of the Los Catutos Member (Vaca Muerta Formation), Sierra de la Vaca Muerta, Neuquén Basin. \emph{Journal of South American Earth Sciences}, 120, 104081.
  \item \textbf{Likerman, J.}, Zlotnik, S., \& Li, C.-F. (2021). The effects of small-scale convection in the shallow lithosphere of the North Atlantic. \emph{Geophysical Journal International}, 227(3), 1512--1522.
  \item García Morabito, E., Beltrán-Triviño, A., Terrizzano, C.M., Bechis, F., \textbf{Likerman, J.}, Von Quadt, A., \& Ramos, V.A. (2021). The influence of climate on the dynamics of mountain building within the Northern Patagonian Andes. \emph{Tectonics}, 40(2), e2020TC006374.
  \item Luna, C.C., Yagupsky, D.L., \& \textbf{Likerman, J.} (2019). Fracture network analysis of Yacoraite Formation in the Tres Cruces sub-basin, northwestern Argentina. \emph{Journal of South American Earth Sciences}, 102446.
  \item Jara, P., \textbf{Likerman, J.}, Charrier, R., Herrera, S., Pinto, L., Villarroel, M., \& Winocur, D. (2017). Closure type effects on the structural pattern of an inverted extensional basin of variable width: analogue models. \emph{Journal of South American Earth Sciences}.
  \item Gianni, G.M., Echaurren, A., Folguera, A., \textbf{Likerman, J.}, Encinas, A., García, H.P.A., Dal Molin, C., \& Valencia, V.A. (2017). Cenozoic intraplate tectonics in Central Patagonia: Record of main Andean phases in a weak upper plate. \emph{Tectonophysics}, 721, 151--166.
  \item Terrizzano, C.M., García Morabito, E., Christl, M., \textbf{Likerman, J.}, Tobal, J., Yamin, M., \& Zech, R. (2017). Climatic and tectonic forcing on alluvial fans in the Southern Central Andes. \emph{Quaternary Science Reviews}, 172, 131--141.
  \item Cilona, A., Aydin, A., \textbf{Likerman, J.}, Parker, B., \& Cherry, J. (2016). Structural and statistical characterization of joints and multi-scale faults in an alternating sandstone and shale turbidite sequence at the Santa Susana Field Laboratory. \emph{Journal of Structural Geology}, 85, 95--114.
  \item Tobal, J.E., Folguera, A., \textbf{Likerman, J.}, Naipauer, M., Sellés, D., Boedo, F.L., Ramos, V.A., \& Gimenez, M. (2015). Middle to late Miocene extensional collapse of the North Patagonian Andes (41°30'--42°\,S). \emph{Tectonophysics}, 657, 155--171.
  \item Ghiglione, M.C., \textbf{Likerman, J.}, Barberón, V., Giambiagi, L.B., Aguirre-Urreta, B., \& Suarez, F. (2014). Geodynamic context for the deposition of coarse-grained deep-water axial channel systems in the Patagonian Andes. \emph{Basin Research}, 26(6), 726--745.
  \item Jara, P., \textbf{Likerman, J.}, Winocur, D., Ghiglione, M.C., Cristallini, E.O., Pinto, L., \& Charrier, R. (2014). Role of basin width variation in tectonic inversion: implications for the Abanico Basin, 32°--34°\,S, Central Andes. \emph{Geological Society, London, Special Publications}, 399, SP399-7.
  \item \textbf{Likerman, J.}, Burlando, J.F., Cristallini, E.O., \& Ghiglione, M.C. (2013). Along-strike structural variations in the Southern Patagonian Andes: Insights from physical modeling. \emph{Tectonophysics}, 590, 106--120.
\end{enumerate}
\end{multicols}
}

%-------------------------------------------------------------------------------
\cvsection{Fellowships \& Distinctions}
\begin{cvhonors}
\cvhonor{Research Fellowship in Japan}{Kobe University}{JSPS}{2027}
\cvhonor{SZNet 2025 Chile Field Trip}{National Science Foundation}{SZ4D}{2025}
\cvhonor{Academic Mobility Fellowship among Partner Institutions}{Ibero-American University Association for Postgraduate Studies}{AUIP}{2022}
\cvhonor{International Research Stay}{National Scientific and Technical Research Council}{CONICET}{2016}
\cvhonor{Postdoctoral Fellowship}{National Scientific and Technical Research Council}{CONICET}{2015}
\cvhonor{Doctoral Research Stay at Stanford University}{Fulbright -- Bunge \& Born Foundation}{Doctoral fellowship}{2013}
\end{cvhonors}

%-------------------------------------------------------------------------------
\cvsection{Skills \& Languages}
\descriptionstyle{\textbf{Modelling:} Underworld2, LaMEM, GOLEM, ASPECT; analogue modelling, geomechanics, finite-element modelling; HPC workflows.}\par
\vspace{0.8mm}
\descriptionstyle{\textbf{Spatial analysis \& fieldwork:} QGIS, GMT, PostGIS; seismic interpretation, geological mapping, UAV-assisted methods; 15+ years leading Andean field campaigns.}\par
\vspace{0.8mm}
\descriptionstyle{\textbf{Programming:} Python, R, MATLAB, \LaTeX{}, HTML5.}\par
\vspace{0.8mm}
\descriptionstyle{\textbf{Languages:} Spanish (native), English (fluent, C2), Italian (A2), Catalan (A2).}

%-------------------------------------------------------------------------------
\cvsection{Alignment with Call Priorities}
\begin{cvparagraph}
\textbf{Priority themes:} direct, ongoing research programme in Geodynamics (subduction dynamics, mantle-plume interactions, lithospheric deformation) and Earth Resources (natural-fracture prediction in unconventional reservoirs; geothermal-potential assessment), evidenced by a Principal-Investigator-led project track record across five internationally co-funded grants.

\textbf{Capacity-building:} continuous university teaching (2013--present), a dedicated professional-training course and field-based instruction for early-career geologists, and supervision of one postdoctoral researcher, one doctoral thesis, and eleven undergraduate theses.

\textbf{Regional cooperation:} documented, multi-year collaboration with Universidad de Chile researchers producing peer-reviewed outputs within the same publication series as a prior IGCP project (586-Y), continued engagement with the NSF SZ4D regional subduction-zone network, and current management of a bi-national (Argentina--Spain) research cooperation, providing direct experience relevant to leading a binational or regional IGCP project across Latin America and the Caribbean.
\end{cvparagraph}

\end{document}
